% ED26-SRC-027-501: source section3 PA/Puel continuation; not yet book-routed.
\subsection{Norm-indexed PA: the contextual obligation and its repair}
\label{sec:p12-klop1987-pa-norm-proof}

Remark3.11(ii), printed p.21, gives the three ordinary rules
\begin{align*}
 x\parallel y&\to(x\mathbin{\lhd}y)+(y\mathbin{\lhd}x),\\
 (x+y)\mathbin{\lhd}z&\to(x\mathbin{\lhd}z)+(y\mathbin{\lhd}z),\\
 (a\cdot x)\mathbin{\lhd}y&\to a\cdot(x\parallel y).
\end{align*}
Here \(\lhd\) denotes the source's binary left-merge glyph, not
an inequality; \(a\) ranges over atomic-action constants. The
unindexed precedence \(\parallel>\lhd>\cdot,+\) does not handle
the last rule, which recreates the heavier operator. The report
introduces indexed operators and illustrates the method on its
node-count-preserving \(g/h/f\) pair. That pair's proof does not
by itself justify annotations above a duplicating PA step
\parencite[p.21]{Klop1987CWITutorial}.

The earlier primary CWI report CS-R8403 has the same indexing idea
in Appendix I,13, but defines an additive atomic-length norm, including
addition across \(+\), rather than counting every syntax node
\parencite[pp.30--34]{BergstraKlop1984PAIndexingReport}.
These are different norms. The construction below both exhibits the
contextual reannotation failure and supplies a complete repaired
max/sum interpretation. It also gives a budget interpretation that
can retain the tutorial's original \emph{initial} syntax-length tags.
Neither repair is silently identified with a full proof printed in
either source.

\begin{hbdefinition}{def:p12-klop1987-pa-regime}{The exact ordinary PA rewrite systems being proved}
Let \(A\) be a nonempty set of distinct nullary atomic-action symbols
and \(V\) an infinite variable supply disjoint from it. Finite free
terms have those leaves and four binary symbols
\(+,\cdot,\parallel,\lhd\). Let \(R_3\) consist exactly of
the three displayed rules, including the action-indexed third schema.
Let \(R_7\) additionally contain
\[
 a\mathbin{\lhd}x\to a\cdot x,\quad x+x\to x,\quad
 (x+y)\cdot z\to(x\cdot z)+(y\cdot z),\quad
 (x\cdot y)\cdot z\to x\cdot(y\cdot z).
\]
The first additional schema is the atomic left-merge rule; the other
three are the oriented idempotence, distribution and sequential
associativity equations. The primary table calls these CM2, A3,
A4 and A5. Its ACP merge equation also includes a communication
merge summand, which is \emph{not} in the tutorial's PA rule above.
That primary report assumes a finite action alphabet; the proof below
also covers an infinite ambient action set by establishing a finite
alphabet for each individual run, rather than attributing that
extension to the primary statement.

All rewriting here is ordinary rewriting on these free finite terms,
under arbitrary finite-term instances and contexts. Commutativity
and associativity of \(+\) are not directed rules or implicit
matching equations. No communication merge, deadlock, silent action,
encapsulation, abstraction or recursive-definition rule is added.
Thus \(R_7\) is the expressly stated ordinary PA presentation;
it is not all of the primary report's larger ACP system. That
appendix in fact excludes A5 from its commutative-tree argument;
the present A5 case uses ordered left-to-right LPO separately.
\end{hbdefinition}

\begin{hblemma}{thm:p12-klop1987-pa-reannotation-failure}{proved=annex:p12-klop1987-pa-reannotation-failure}{Literal length reannotation does not lift every contextual PA step}
Take \(t=(a\parallel b)\mathbin{\lhd}c\) and perform the first
source rule inside its left argument, giving
\(u=((a\mathbin{\lhd}b)+(b\mathbin{\lhd}a))\mathbin{\lhd}c\).
Annotating every merge in every successive term afresh by the sum
of its arguments' syntax-node counts changes the outer left-merge
index from4 to8. With the primary additive atomic-length norm applied
to this two-branch PA rule it changes from3 to5. Applying the primary
report's actual three-branch CM1 rule, with its additional communication
merge \(a\mathbin{|}b\), changes that primary norm's outer index
from3 to7. In each case the new outer numerical label
exceeds every old numerical label in any order realizing the source's
interleaved indexed precedence. The source marker operations cannot
produce such a larger label. Therefore this automatic reannotation
does not give the required positive marker comparison in an arbitrary
context. This refutes that proposed lift, not termination of PA.
\end{hblemma}

\begin{hbdefinition}{def:p12-klop1987-pa-stable-norm}{A positive max/sum norm and its canonical annotation}
Define a different norm \(\nu\) by
\[
 \nu(a)=\nu(x)=1,\qquad
 \nu(p+q)=\max(\nu(p),\nu(q)),\qquad
 \nu(p\circ q)=\nu(p)+\nu(q)
 \quad(\circ\in\{\cdot,\parallel,\lhd\}).
\]
Thus every action and variable has positive unit weight, and choice
uses maximum, not sum or syntax size. Let \(E(t)\) preserve variables,
actions, \(+\) and \(\cdot\), recursively annotate the children,
and replace \(p\parallel q\) and \(p\lhd q\) by
\(P_{\nu(p)+\nu(q)}(E(p),E(q))\) and
\(L_{\nu(p)+\nu(q)}(E(p),E(q))\), respectively.
All \(P_n,L_n\) have fixed arity two and \(n\geq2\).

Use left-to-right LPO with precedence ordered first by indexed level,
then \(P_n>L_n\) within a level, and all indexed symbols above
\(\cdot>+>a\) for every \(a\in A\). No comparison between
different actions is needed. In particular
\(P_n>L_n\) and \(L_{n+1}>P_n\), retaining the source's
required interleaving. This is an infinite ambient indexed signature;
the run-level finite-signature obligation is proved below.
\end{hbdefinition}

\begin{hblemma}{thm:p12-klop1987-pa-norm}{proved=annex:p12-klop1987-pa-norm}{Invariance for every instance and every ancestor}
The norm is a well-defined positive integer on every finite term;
every subterm has norm at most that of the whole term. Every
\(R_7\)-step preserves it, including under arbitrary substitutions
and contexts. Norm evaluation after a substitution assigns each
variable the positive norm of its substituted term; no unjustified
identity between an indexed annotation and an ordinary substitution
is assumed. A change with equal norm leaves all ancestor merge
indices unchanged, so both endpoints have one common annotated
one-hole context around the changed subtree.
\end{hblemma}

For example, \(E(x\parallel y)=P_2(x,y)\), but substituting
\(a\cdot a\) for \(x\) gives canonical annotation
\(P_3(a\cdot a,y)\), not \(P_2(a\cdot a,y)\).
Consequently the proof checks actual instances first, rather than
asserting that ordinary substitution preserves numerical annotations.

\begin{hblemma}{thm:p12-klop1987-pa-lpo-lift}{proved=annex:p12-klop1987-pa-lpo-lift}{Every actual PA rule instance strictly decreases LPO}
For every ordinary instance \(l\sigma\to r\sigma\) of each
of the seven stated rules,
\(E(l\sigma)>_{\rm lpo}E(r\sigma)\). Consequently every
contextual step \(t\to_{R_7}u\) gives
\(E(t)>_{\rm lpo}E(u)\). The proof supplies all domination
and equal-head lexicographic cases, including both independently
sized branches of distribution and the recreated
\(P_{n-1}\) below a consuming \(L_n\).
\end{hblemma}

\begin{hbtheorem}{thm:p12-klop1987-pa-termination}{proved=annex:p12-klop1987-pa-termination}{The complete ordinary PA termination argument}
Both \(R_3\) and \(R_7\) terminate on all finite open and ground
terms of their declared signatures. In a run starting at \(t\),
every annotated merge index is between2 and \(N=\nu(t)\).
Only the finitely many action and variable symbols already in \(t\),
the two fixed base operators, and those finitely many indexed operators
are needed throughout that entire run. There is no common global
bound on indices or term size imposed on all inputs.
The finite-ranked LPO well-foundedness proof therefore applies to
the actual run, rather than being invoked over an unverified infinite
signature. The result is SN, not a confluence or PA completeness theorem.
\end{hbtheorem}

\begin{hbdefinition}{def:p12-klop1987-pa-budget}{Retained initial tags as upper budgets}
For the three-rule system \(R_3\), an admissible annotated term
has erasure in the ordinary signature and satisfies
\(n\geq\nu(p)+\nu(q)\) at each \(P_n(p,q)\) or
\(L_n(p,q)\), with \(\nu\) evaluated on erasures.
This is an upper-bound condition, not equality to the current
syntax length. Its lifted rules are
\begin{align*}
 P_n(p,q)&\to L_n(p,q)+L_n(q,p),\\
 L_n(p+q,r)&\to L_n(p,r)+L_n(q,r),\\
 L_n(a\cdot p,q)&\to a\cdot P_{n-1}(p,q).
\end{align*}
They keep all surrounding tags and use the inherited \(n\) in
each distributed branch. The last rule is defined only where
admissibility gives \(n\geq3\).
\end{hbdefinition}

\begin{hbtheorem}{thm:p12-klop1987-pa-budget-lift}{proved=annex:p12-klop1987-pa-budget-lift}{A complete lift retaining the tutorial's initial syntax-length indices}
Every finite ordinary term has an admissible annotation obtained by
the tutorial's initial syntax-length prescription. From any admissible
annotation of an \(R_3\)-source, every actual ordinary step lifts
to the displayed annotated rule in the same unchanged-tag context.
Admissibility persists, and each lifted step strictly decreases the
stated LPO. If an initial index is present, newly created indices
do not exceed the largest initial index; without any merge node,
the term has no \(R_3\)-step. This supplies another complete termination route for the
actual three source rules, without falsely reannotating every later
term by its enlarged syntax length. It does not claim a coherent
upper-budget lift for every nonlinear or larger ACP rule.
\end{hbtheorem}

\subsection{The distinct embedded-left-side example}

The report's Remark3.11(i) refers the obstruction
\(f(s(x))\to g(s(x),f(p(s(x))))\) to Puel's extension of
Kruskal's theorem. Its left side embeds in its right side by mapping
the first \(f\) to the new inner \(f\), and diving through
the inserted \(p\) before matching \(s\) and its descendants.
Thus a simplification order containing reverse embedding cannot
orient it, as established in
Corollary~\ref{thm:p12-klop1987-mt-obstruction}. That obstruction
is different from PA's re-emerging heavy operator and is not removed
merely by renaming its symbols as merge operators.

\begin{hbtheorem}{thm:p12-klop1987-pa-puel-example}{proved=annex:p12-klop1987-pa-puel-example}{Termination of the actual isolated source example by a different measure}
For the one-rule ordinary TRS
\[
 Q=\{f(s(x))\to g(s(x),f(p(s(x))))\},
\]
with unary \(f,s,p\), binary \(g\), variables and any inert
constants, every finite term terminates. No rule for \(p\) or
\(s\) is present. The proof traces the initially active \(f\)
occurrences: a contracted occurrence disappears, the new \(f\)
above a \(p\)-root is permanently inactive, and only strictly
deeper old active origins are copied. A finite lexicographic residual
count proves termination despite the growth of syntax.
This is a concrete proof of the isolated displayed system, not a
reconstruction of Puel's general unavoidable-set theorem.
\end{hbtheorem}

\begin{hblemma}{thm:p12-klop1987-pa-puel-extension-failure}{proved=annex:p12-klop1987-pa-puel-extension-failure}{Why the isolated proof does not settle arbitrary companion rules}
Adding \(p(s(x))\to s(x)\) destroys termination. For an inert
constant \(c\) (or a fixed object variable), two steps take \(f(s(c))\) to
\(g(s(c),f(s(c)))\); repetition in the new second-child context
gives an infinite ordinary reduction. The absence of companion
\(p\)-rules in the preceding theorem is therefore substantive.
\end{hblemma}

\begin{proofboundary}{def:p12-klop1987-pa-boundaries}{Separate norms, separate equations, separate tree extensions}
The original length annotations, the repaired max/sum equality
annotations and the inherited upper budgets have distinct meanings.
The proofs above establish the exact three PA rules and the expressly
listed seven-rule ordinary presentation, not all ACP rules or all
rewriting modulo commutative/associative choice. Directing choice
commutativity would itself permit an infinite swapping reduction.
The primary CWI appendix's indexed communication, silent-action and
abstraction cases require their own norm and lifting proofs.
Puel's cited report is a different tree-order extension; neither
the PA proof nor the isolated active-origin argument establishes its
general well-quasi-order hypotheses or its unavoidable-set construction.
The distinct 1987 and 1992 presentations retain their own source standing.
\end{proofboundary}
